\section{Learning a sparse support from a short row}
\label{learn:section}

The remaining low-dimensional description problem is the following.
A set $S$ has unusually few incidences with a large opposite set $T$.
We want a short message describing a set that captures a fixed fraction
of $S$ and is not much larger than $S$. The decoder does not know either
support. We first construct the desired set from a random row of points
of a prepared subset of $S$, and then implement that row by searching a
public table of proposals from the known domain.

Two estimates have different roles. A second moment will show that the
test retains most points of the hidden support. A higher moment will
exclude almost all unwanted ambient points. Keeping these roles separate
is essential: some geometric exceptions allowed in the ambient estimate
are too large to discard from the support itself.

\subsection{The prepared input and the sampling test}

Use the parameters of Equation~\eqref{base:parameters}, with
$\sigma=\log q$, and put
\begin{equation}\label{learn:auxiliary-parameters}
 \chi=P/100,\qquad \xi=L/100,\qquad
 p\in[10000\sigma/P,10000\sigma/P+2]\cap2\mathbb Z,
 \qquad t_*=(100p)^{-1}.
\end{equation}
Here $p$ is a moment order, not the characteristic; the field is $\F_q$
with $q$ prime. In particular,
\begin{equation}\label{learn:scale-use}
 b+P\tau=o(L),\quad R\log L=o(P),\quad
 \log\sigma=o(L),\quad R\gg\log(p+2),\quad P=o(\sigma).
\end{equation}
These relations follow from Lemma~\ref{base:scale-comparisons} whenever
$b=O(K_*)$ and $\tau=O(\sigma^{-100\beta})$.

We specify the geometric input to this section. Let $j\in\{2,3\}$,
let $S,T$ be nonempty sets in opposite projective spaces of dimension
$j$, and suppose that
\begin{equation}\label{learn:support-hypotheses}
 S\subseteq U_S,\qquad n=|S|\le |T|,\qquad
 |S||T|\ge q^{j+1}e^{-b},\qquad
 e(S,T)\le \frac{\tau}{q}|S||T|.
\end{equation}
The set $U_S$ is known to the decoder. Write
$d_S=\log(|U_S|/n)$. We only need to treat
$n>100qP$, since smaller supports can be enumerated.

The prepared support is a subset $S'\subseteq S$ of size
$n'\in[n/2,n]$. In dimension three there may also be an ordered list
of projective planes $\Pi_1,\ldots,\Pi_J$. Its cells are
\[
 C_i=S'\cap\bigl(\Pi_i\setminus\textstyle\bigcup_{a<i}\Pi_a\bigr).
\]
The list is allowed to be empty; in dimension two it is always empty.
The list and the sizes of its cells admit an $O(q)$-bit description, and
every cell has size at most $.04n'$. For an ambient point $x$, let
$O_x=C_i$ if $\Pi_i$ is the first listed plane containing $x$, and let
$O_x=\varnothing$ if there is no such plane. Set
\begin{equation}\label{learn:radial-definitions}
 f_x=|O_x|/n',\quad X_x=S'\setminus(O_x\cup\{x\}),\quad
 B=q^j/n',\quad n=q^{1+v}.
\end{equation}
Thus $0\le f_x\le .04$. Every line through $x$ is called a radial
line at $x$, and has strength
$a_\ell=q|\ell\cap X_x|/n'$.

For each $x$, consider any family of hyperplanes $H$ through $x$ with
$q|H\cap X_x|/n'\le2$. For two distinct members put
$\theta_{HH'}=q|H\cap H'\cap X_x|/n'$. Let $E_a$ be the number of
ordered distinct pairs with $\theta_{HH'}\ge a$, and let $\Delta_a$
be the largest number of such neighbors of one member. The geometric
preparation supplies a deterministic exceptional set $D_0$ of size
$o(n)$ such that, outside $D_0$, simultaneously for dyadic
$a\in[q^{-10},2]$,
\begin{equation}\label{learn:overlap-input}
 E_a\le B^2e^{2\chi}a^{-200},\qquad
 \Delta_a\le Be^{2\chi}a^{-200}.
\end{equation}
There is a further deterministic set $D_1$, of size at most
$n e^{C_0P}$, outside which, also excluding $D_0$,
\begin{equation}\label{learn:strong-degree-input}
 \Delta_{t_*}\le Be^{-3P}.
\end{equation}
The bounds apply to all the indicated families. In dimension two the
overlaps vanish, so both assertions hold without geometric exceptions.
Proposition~\ref{geom:preparation} establishes this prepared input after its
enumeration and large-cell reductions, with the overlap bounds supplied
by Lemma~\ref{geom:overlap}.

\begin{proposition}[Short-row description]\label{learn:short-row}
Fix the constants in \eqref{learn:support-hypotheses}, the bound on
$b/K_*$, the bound on $\tau\sigma^{100\beta}$, and $C_0$.
For every sufficiently large prime $q$, the prepared input just
specified has the following description scheme. Using a fresh public
table, it either declares failure or describes a set $W\subseteq U_S$
satisfying
\[
 |W|\le n e^{CP},\qquad |W\cap S|\ge c n.
\]
Failure has probability at most $C e^{-cq}$, and the message has at most
$CqP(d_S+1)+Cq+C\sigma$ bits. In particular this is within
$CqP(d_S+d_T+P)$ for any $d_T\ge0$.
The bounds are uniform in all supports and prepared inputs satisfying
the hypotheses, including inputs fixed by earlier random history before
the fresh table is sampled. The decoder uses only $U_S$, public
parameters, and the new message; it need not know $S$, $T$, or $S'$.
\end{proposition}

Mixing in Lemma~\ref{base:incidence} and the smaller-side orientation
give $n\le Cq^{(j+1)/2}$. Consequently the quantities used below obey
\begin{equation}\label{learn:elementary-sizes}
 \frac q{n'}<\frac1{50P},\qquad
 Q_{j-1}=O(B^2),\qquad B^{-1}\le e^{-P},\qquad
 v\le1+O(1/\sigma).
\end{equation}
The last bound will only be needed in dimension three.

We first sample from the hidden set $S'$. Take $R$ independent batches.
Each batch has Poisson size of mean $qL$, and, conditional on its size,
has independent uniform entries in $S'$. Equivalently, the numbers of
occurrences of the different point-batch pairs are independent Poisson
variables of mean $qL/n'$. Indeed the joint probability of prescribed
counts is obtained by multiplying the Poisson probability of their sum
by the multinomial probability conditional on that sum; the factorials
then factor over points. The total number of samples has mean $qP$.

For a hyperplane $H$ and a cell $C$, define
\[
 \lambda_H=\frac{q|S'\cap H|}{n'},\qquad
 \lambda_{H,C}=\frac{q|C\cap H|}{n'},\qquad f_C=\frac{|C|}{n'}.
\]
Call $H$ exceptional if $|\lambda_H-1|>.1$ or
$|\lambda_{H,C}-f_C|>.02$ for at least one cell; let $\mathcal E$
be this set. All other hyperplanes are typical. Let
\begin{equation}\label{learn:empty-set}
 Z=\{H\in\mathcal E:H\text{ has no sample in any batch}\},\qquad
 z=|Z|,\qquad z_x=|\{H\in Z:x\in H\}|.
\end{equation}
These objects are used by the encoder; the decoder will receive only
the integer $z$, not the hidden set $Z$.

For a query point $x$, put $b_x=e^{-L(1-f_x)}$, and let $N_r(A)$
denote the number of samples of batch $r$ in a point set $A$.
For each hyperplane through $x$, we multiply $R$ batch tests.
In each batch we center the emptiness indicator outside $O_x$ at $b_x$,
and set the test to zero if $O_x\cap H$ contains a sample. Thus a
hyperplane empty in every batch contributes $(1-b_x)^R$, whereas
any other hyperplane has product of absolute value at most $b_x$.
For typical hyperplanes, their normalized mass outside $O_x$ is close
to $1-f_x$.
Batch independence turns bounds on one-batch means and pair moments
into $R$th-power bounds. The overlap estimates then control the sum
over typical hyperplanes.
Define the score by
\begin{equation}\label{learn:score}
 A_x=\sum_{H\ni x}\prod_{r=1}^R
 \left[
 1_{\{N_r(O_x\cap H)=0\}}
 \left(1_{\{N_r((S'\setminus O_x)\cap H)=0\}}-b_x\right)
 \right].
\end{equation}
The random test includes every sampled point and every other point
with $A_x<z/(2q)$. Both factors inside the brackets belong to batch
$r$. Although the formula mentions hidden subsets, the decoder can
evaluate it from the samples, the plane list, the cell sizes, and $z$:
earliest membership in the list identifies each sample's cell and the
cell $O_x$, if present. The size $n'$ is transmitted as well.
The separate own-cell indicator also preserves the geometric benefit
of removing $O_x$: once its sample counts and the absence of a sample
at $x$ are fixed, the remaining randomness consists of independent
counts on the radial lines in $X_x$. We use this conditioning when
estimating the typical contribution.

\subsection{Regular pencils and empty exceptional hyperplanes}

The pencil at $x$ is the set of hyperplanes through $x$. We next discard
a small deterministic set of pencils on which squared deviations are
unusually concentrated. This is separate from all sampling events.

\begin{lemma}[Regular pencils]\label{learn:regular-pencils}
There is a deterministic set $D_2$ of size $o(n)$ such that, for
$x\notin D_2$, its pencil contains at most $Be^\xi$ exceptional
hyperplanes and satisfies
\[
 \sum_{H\ni x}\min\{1,(\lambda_H-1)^2\}\le Be^\xi,
 \qquad
 \sum_{H\ni x}\min\{1,(\lambda_{H,C}-f_C)^2\}\le Be^\xi
\]
for every cell $C$. The constants and the $o(n)$ estimate are uniform
over the prepared inputs. We call these pencils regular.
Moreover, for every regular pencil and its typical members, with
\begin{equation}\label{learn:outside-mass}
 \lambda'_H=\frac{q|H\cap X_x|}{n'},\qquad
 \delta'_H=\lambda'_H-(1-f_x),
\end{equation}
we have
\begin{equation}\label{learn:typical-bounds}
 .75\le\lambda'_H<2,\qquad |\delta'_H|\le .18,\qquad
 \sum_{\substack{H\ni x\\H\text{ typical}}}(\delta'_H)^2
 \le CBe^\xi.
\end{equation}
\end{lemma}

\begin{proof}
Weighted incidence variance gives
\begin{equation}\label{learn:total-variance}
 \sum_H(\lambda_H-1)^2\le CqB,\qquad
 \sum_H(\lambda_{H,C}-f_C)^2\le CqBf_C.
\end{equation}
The exact means differ from $1$ and $f_C$ by the factors
$qp_j=1-1/Q_j$, whose squared contributions satisfy the same bounds.
Summing over cells and using $\sum_C f_C\le1$ proves
\begin{equation}\label{learn:exceptional-size}
 |\mathcal E|\le CqB,\qquad
 \sum_{H\in\mathcal E}\lambda_H\le CqB.
\end{equation}
For the second assertion use Cauchy--Schwarz on
$\sum_{H\in\mathcal E}(\lambda_H-1)$.

For the clipped squared-deviation weight of a cell, the total weight is
at most $CqBf_C$ and each weight is at most one. Its pencil sums have
mean $O(Bf_C)$ and total squared deviation at most $Cq^jBf_C$, by
Lemma~\ref{base:incidence}. For sufficiently large $q$, exceeding
$Be^\xi$ therefore excludes at most
\[
 C\frac{q^jBf_C}{B^2e^{2\xi}}=Cn'f_Ce^{-2\xi}
\]
centers. Sum this estimate over cells, rather than multiplying by their
number. The full-deviation weight and the exceptional-indicator weight
have the same proof with fraction one. Their union is $D_2$.

For a typical hyperplane, removing $O_x$ subtracts a quantity within
$.02$ of $f_x$, when a cell is present. Removing $x$ subtracts at most
$q/n'$. Thus $\lambda'_H\ge .9-.04-.02-q/n'>.75$ and
$|\delta'_H|\le .1+.02+q/n'<.18$. The upper bound by 2 follows as well.
On typical hyperplanes the relevant squared deviations are already
below one. The inequality $(u+v+w)^2\le3(u^2+v^2+w^2)$ and regularity
give the last bound in \eqref{learn:typical-bounds}; the removed center
adds at most
$O(q^{j-1}(q/n')^2)=O(B)$ by \eqref{learn:elementary-sizes}.
\end{proof}

\begin{lemma}[The empty exceptional set]\label{learn:empty-exceptional}
There is a fixed $c_0>0$ such that, with probability at least
$c e^{-8P\tau}$,
\begin{equation}\label{learn:empty-good-event}
 z\ge z_{\min}:=c_0qB e^{-b-8P\tau},\qquad
 \sum_{H\in Z}\lambda_H\le .001z.
\end{equation}
For every $Z$ satisfying these two inequalities, at least $.99n'$
points of $S'$ have $z_x\le .1z/q$, whereas all but
$O(q^{j+1}/z)$ ambient points have $z_x\ge .8z/q$.
At a regular pencil, the exceptional-hyperplane part of
\eqref{learn:score} differs from $(1-b_x)^Rz_x$ by at most
$b_xBe^\xi=o(z_{\min}/q)$, uniformly in the samples, and
$(1-b_x)^R=1-o(1)$.
\end{lemma}

\begin{proof}
The average of $\lambda_H$ over $T$ is at most $2\tau$.
At least half of $T$ consequently has $\lambda_H\le8\tau$.
These hyperplanes are exceptional, and their number $m$ is at least
$c qB e^{-b}$ by \eqref{learn:support-hypotheses} and $n'\ge n/2$.
Each is empty with probability at least $a=e^{-8P\tau}$.
If $Y$ counts the empty members of this particular set, then
$0\le Y\le m$ and $\E Y\ge ma$. Therefore
$\Prob(Y\ge ma/2)\ge a/2$, without any independence assumption among
the hyperplanes. Choose $c_0$ small enough to obtain the first inequality
of \eqref{learn:empty-good-event} on this event.

By \eqref{learn:exceptional-size}, the expected contribution from
$\lambda_H\ge10^{-4}$ to $\sum_{H\in Z}\lambda_H$ is at most
$CqB e^{-10^{-4}P}$. Markov's inequality at $.0009z_{\min}$ bounds
the probability of a larger contribution by
\[
 C\exp\{b+8P\tau-10^{-4}P\}=o(e^{-8P\tau}).
\]
The remaining contribution is at most $.0001z$. Subtracting this
failure probability proves the asserted simultaneous event.

For fixed $Z$, incidence variance has mean $p_jz=(1-o(1))z/q$ and
total squared deviation at most $q^{j-1}z$. This gives the ambient
exception bound. On the hidden support, exact double counting gives
\[
 \sum_{x\in S'}z_x=\frac{n'}q\sum_{H\in Z}\lambda_H
 \le .001n'z/q,
\]
which proves the $.99n'$ assertion by Markov's inequality.

An exceptional hyperplane in $Z$ contributes $(1-b_x)^R$ to the score.
For any other exceptional member of the pencil, an own-cell hit makes
its product zero; otherwise an outside hit gives a factor $-b_x$,
and all other factors have absolute value at most one. Its contribution
has absolute value at most $b_x$. Regularity proves the error bound.
Finally $b_x\le e^{-.96L}$, so
$b_xBe^\xi\le Be^{-.95L}=o(z_{\min}/q)$ and $Rb_x=o(1)$ by
\eqref{learn:scale-use}.
\end{proof}

The empty hyperplanes thus give a score that is small on most of $S'$
and large on most ambient points. It remains to show that the typical
hyperplanes do not obscure this distinction. Write $T_x^*$ for their
contribution to \eqref{learn:score}.

\subsection{Conditional moments of the typical contribution}

Fix $x$ and condition on its being unsampled and on all point-batch
counts in $O_x\setminus\{x\}$. Denote these conditioning data by
$\mathcal O_x$; when $x\notin S'$, the unsampled condition is automatic.
Every assertion below is uniform over conditioning values of positive
probability. The own-cell multipliers in \eqref{learn:score} are now
fixed zeros or ones. Counts on $X_x$ are unchanged independent Poisson
variables. Grouping them by radial lines gives independent
line-batch schedules, with hit probability at most $La_\ell$ for
line $\ell$ in one batch. Distinct lines are disjoint in $X_x$ because
the center has been removed.

Temporarily omit a fixed own-cell multiplier and let
$X_H=1_{\{H\cap X_x\text{ empty in the batch}\}}-b_x$.
Its mean $u_H$ and its pair moments obey
\begin{align}
 u_H&=e^{-L\lambda'_H}-b_x,
 &|u_H|&\le Le^{-.75L}|\delta'_H|,\label{learn:one-mean}\\
 \E(X_HX_{H'})
 &=u_Hu_{H'}+e^{-L(\lambda'_H+\lambda'_{H'})}
                      (e^{L\theta_{HH'}}-1),\label{learn:one-covariance}\\
 |\E(X_HX_{H'})|
 &\le L^2e^{-.75L}(|\delta'_H\delta'_{H'}|+\theta_{HH'}),
 &\E X_H^2&\le2e^{-.75L}.\label{learn:one-bounds}
\end{align}
The covariance identity is the Poisson emptiness formula for a union.
For its bound, use
$\lambda'_H+\lambda'_{H'}-\theta_{HH'}\ge .75$ and
$e^u-1\le ue^u$ for $u\ge0$.

\begin{lemma}[Second moment]\label{learn:second-moment}
If $x\notin D_0\cup D_2$, then, conditional on $\mathcal O_x$,
\[
 \E\bigl((T_x^*)^2\mid\mathcal O_x\bigr)
 \le B^2e^{-.6P}.
\]
This assertion does not exclude the set $D_1$.
\end{lemma}

\begin{proof}
Batch independence raises the bounds in \eqref{learn:one-bounds} to
the $R$th power. Fixed zero-one multipliers cannot increase their absolute
values. By \eqref{learn:typical-bounds},
\[
 \sum_H|\delta'_H|^R\le CBe^\xi.
\]
For the overlap terms, \eqref{learn:overlap-input} and dyadic summation
give, for $R>200$,
\[
 \sum_{H\ne H'}\theta_{HH'}^R
 \le B^2\exp\{2\chi+O(R+\log\sigma)\}.
\]
Every positive overlap is at least $q/n'$ and at most 2, so it falls
within the available dyadic range. There are $O(\sigma)$ dyads. The
inequality $(u+v)^R\le2^{R-1}(u^R+v^R)$ handles their sum with the
deviation terms. The diagonal contributes at most
$Q_{j-1}(2e^{-.75L})^R$. Now use $Q_{j-1}=O(B^2)$,
$2\chi=.02P$, and $R\log L=o(P)$ to obtain the claimed exponent.
\end{proof}

\begin{lemma}[High moment]\label{learn:high-moment}
If $x\notin D_0\cup D_1\cup D_2$, then, conditional on
$\mathcal O_x$,
\[
 \E\bigl((T_x^*)^p\mid\mathcal O_x\bigr)
 \le B^p e^{-.1pP}.
\]
\end{lemma}

\begin{proof}
We expand into ordered $p$-tuples of typical hyperplanes and bound the
absolute expectation of every summand. The evenness of $p$ will permit
the final tail estimate. A tuple may contain repeated hyperplanes.
Among its distinct hyperplanes, join $H,H'$ when
$\theta_{HH'}>t_*$, and take the connected components of this graph.
A component consisting of one hyperplane appearing once in the tuple
is a \emph{simple singleton}. All other components are nonsimple.

\paragraph{Component expectations.}
A positive-strength radial line is external to the tuple if it is
contained in hyperplanes in different components. The sum of the
strengths of the external lines in any fixed tuple hyperplane is at
most $pt_*=.01$: assign each such line to one other tuple hyperplane
containing it, whose overlap is at most $t_*$. Condition further on
all external line-batch schedules. The remaining schedules affecting
different components are independent, and each hyperplane has internal
mass at least $.75-.01=.74$. Call a hyperplane touched in a batch if
one of its external lines has a hit in that batch.

For a nonsimple component, all factor magnitudes are at most one.
Keeping just the factors of one of its hyperplanes bounds its absolute
product expectation. In an untouched batch the expected absolute
factor is at most $2e^{-.74L}$; in a touched batch it is at most
$b_x\le e^{-.96L}$. Independence between batches consequently bounds
the entire component by
\begin{equation}\label{learn:nonsimple-attenuation}
 (2e^{-.74L})^R\le e^{-.7P}.
\end{equation}

For a simple singleton keep its signed expectation. Let
$\theta_{\max}$ be the largest overlap $\theta_{HH'}$ over the other
distinct hyperplanes $H'$ of the current tuple, taking zero if there
are none. It is not a maximum over the whole pencil. In an untouched
batch the mean has absolute value at
most
\[
 Le^{-.74L}(|\delta'_H|+p\theta_{\max});
\]
in every batch it has absolute value at most $2e^{-.74L}$.
Let $\mathcal B_H$ be the event that $H$ is touched in at least
$R/2$ batches. If it is not, at least $R/2$ factors have the sharper
bound. Since $p\theta_{\max}\le .01$, $R/2\ge200$, and
$R\log(2L)=o(P)$, the resulting bound for a simple singleton is
\begin{equation}\label{learn:simple-attenuation}
 e^{-.68P}\left(
 |\delta'_H|^{R/2}+(p\theta_{\max})^{200}+1_{\mathcal B_H}
 \right).
\end{equation}
The fixed own multipliers do not change either estimate.

We now sum these fixed-tuple bounds over the possible hyperplanes.
For each valid tuple, retain $e^{-.7P}$ per nonsimple component and
$e^{-.68P}$ per simple position, and expand the three terms in
\eqref{learn:simple-attenuation} at each simple position. The resulting
bounds are nonnegative. After averaging the external schedules, the
only remaining probabilistic factor in each term is the conditional
probability of the selected events $\mathcal B_H$ occurring together.
We will bound the sum of these terms after dividing by $B^p$.

The geometric choices have three forms. Nonsimple components are
counted either from a strong pair and a spanning tree or from repetitions
of one hyperplane; deviation and overlap terms use the corresponding
pencil bounds; the remaining hit events
are covered by static certificates. Each such certificate charges
any independent direction-batch hit event at most once. We construct
these certificates below and then verify the probability union bound.
These descriptions count moment terms and are not part of the
public-table message.

The combinatorial descriptions use only $p^{O(p)}$ choices of names,
partitions, and orders. Record the partition
of positions into equal hyperplane names, the partition of distinct
names into their strong components, and spanning trees and traversal
orders inside those components. Do not record the full strong graph.
For any geometric assignment its true graph determines whether these
recorded components are valid; an invalid assignment initially contributes
zero. Dropping such validity restrictions later only enlarges a positive
sum after the attenuation bounds have been assigned.

\paragraph{Nonsimple and deviation choices.}
For a component with at least two distinct hyperplanes, start with a
strong pair. By \eqref{learn:overlap-input}, its normalized count is
at most
\[
 e^{2\chi}(C/t_*)^{200}.
\]
Each later distinct vertex along its tree costs at most $e^{-3P}$ by
\eqref{learn:strong-degree-input}; a repeated position costs $1/B$.
For a component consisting of one hyperplane repeated, its first two
positions cost $Q_{j-1}/B^2=O(1)$, and further positions again cost
$1/B$. These counts are combined with
\eqref{learn:nonsimple-attenuation} once per component.

A deviation choice has weighted normalized sum at most
$Ce^\xi$ by \eqref{learn:typical-bounds}. Encode these and all nonsimple
components first. The remaining choices are overlap shifts and events
$\mathcal B_H$; their attenuation $e^{-.68P}$ has already been retained.

\paragraph{Overlap shifts.}
For an unencoded shift vertex, choose the name of a neighbor attaining
its largest positive overlap and a dyad $a\le\theta<2a$. Zero weights
need not be encoded. If the neighbor is encoded, the weighted normalized
count is at most
\[
 (2pa)^{200}\Delta_a/B.
\]
If it is not, encode the ordered pair with count
$(2pa)^{200}E_a/B^2$, dropping the neighbor's extra shift weight or
event requirement, each bounded by one. Both vertices retain their
attenuation. Equation~\eqref{learn:overlap-input} cancels $a^{200}$
in either case and gives at most $e^{2\chi}$ times polynomial overhead.
Summing dyads has the same type of overhead. Continue until all
unencoded vertices carry the event requirement $\mathcal B_H$.

\paragraph{Hit certificates and their probability weights.}
Fix a tuple and external schedules for which these remaining requirements
hold. We describe certificates witnessing enough of the hits to count
the tuple. Maintain two objects: a dictionary of radial directions whose
geometry is already recorded, and a set of charged direction-batch
events. On first introduction of a direction, charge one positive hit
event, except in the pair operation below, which charges 200 distinct
batches at once. Once that direction enters the dictionary, any later
reference to it imposes \emph{no new hit requirement}. In particular,
no direction-batch event is charged twice.

For a fixed tuple all charged events are external events, independent
under the conditional law given $\mathcal O_x$. Thus their joint
probability is at most the product of their $La_\ell$ weights. Keeping
$\min(1,La_\ell)$ would also be valid, but is unnecessary. From an
encoded hyperplane $H_0$, the weighted sum of choices of a new line
and a witnessing batch is at most
\begin{equation}\label{learn:line-charge}
 \sum_{r=1}^R\sum_{\ell\subset H_0}La_\ell
 =P\lambda'_{H_0}<2P.
\end{equation}
An old direction is specified by its dictionary index and requires no
new probability factor. If the current dictionary has $m$ directions,
these references have at most $m$ choices, each of incremental weight
one. A direction choice, allowing either a new charge or an old
reference, therefore has total weight at most $2P+m$. The construction
below keeps $m=O(p)$, so these extra choices enter only the polynomial
overhead.

\paragraph{Anchors and pairs.}
If an unencoded vertex has a touched direction shared with an encoded
hyperplane, encode it using that direction as an anchor. A new direction
is paid for by \eqref{learn:line-charge}; an old one is referenced.
Given the direction, there are at most $Cq^{j-2}$ hyperplanes through
it. The normalized geometric cost is therefore
\[
 O(q^{j-2}/B)=O(q^{v-1})=O(1).
\]
This uses $n'\asymp n$ and \eqref{learn:elementary-sizes}.
In dimension two there are no positive external directions, and no
such event requirements occur.

If no anchor is available but two unencoded vertices share external
hits in at least 200 distinct batches on directions in no encoded
hyperplane, encode this pair. Choose an overlap dyad, 200 witnessing
batches and one shared hit direction in each. These events are new and
distinct, outside the direction dictionary. For a fixed pair their
total probability weights, summed over the witness choices, are at
most $(2Pa)^{200}$. The normalized pair count in
\eqref{learn:overlap-input} then gives a cost at most
$e^{2\chi}(2P)^{200}$. Add the vertices and witness directions to the
encoded data. In dimension three the intersection direction of a
pair is unique; repeated use of that direction in distinct batches
still charges distinct independent events. Repeat the anchor and pair
operations whenever either applies.

\paragraph{The residual root set.}
Suppose $h$ vertices remain. Each has at least $R/2$ touched batches,
whose witnessing external directions are shared with residual vertices
only; otherwise an anchor would be available. Select one witness per
such batch. A fixed neighbor appears fewer than 200 times, and a fixed
direction appears fewer than 200 times. Either violation would permit
the pair operation. Greedily keeping one witness and discarding others
with its neighbor or its direction removes fewer than 400 witnesses.
Thus each residual vertex has at least $\lfloor R/800\rfloor$ witnesses
with distinct neighbors and distinct directions.

There exists a set of at most $h/10$ roots such that every nonroot has
two of these witnesses to roots. To prove this, select each vertex
independently with probability $1/20$. For any one witness list, the
probability of fewer than two selected neighbors is at most
$CR e^{-cR}$. Union over at most $p$ lists is $o(1)$ because
$R\gg\log(p+2)$. Markov's inequality bounds the probability of more
than $h/10$ selected vertices by $1/2$. The two conditions therefore
hold simultaneously with positive probability. If a residual set is
too small to have the displayed number of distinct neighbors, it could
not have reached this case in the first place.

Record the root names and choose their hyperplanes freely, dropping
their extra event requirements. A root has normalized cost
$Q_2/B=O(q^v)$. Encode each nonroot by two distinct directions in
its two chosen root neighbors, using \eqref{learn:line-charge} for
new directions and dictionary references for old ones. Two distinct
radial lines determine at most one projective plane, so the nonroot's
normalized geometric cost is $1/B=O(q^{v-2})$. Restrict to distinct
directions for this uniqueness statement. Each pair operation introduces
at most 200 dictionary directions, each anchor at most one, and each
nonroot at most two. There are at most $p$ tuple vertices, so the
dictionary always has at most $Cp$ entries. Conditional on previously
recorded geometry, each of the two direction choices thus has total
weight at most $2P+Cp$, including old references with unit incremental
weight. After counting the unique plane, their positive weighted sum is
bounded by $(2P+Cp)^2$; this is a polynomial overhead factor.
If there are $r\le h/10$ roots, the aggregate
geometric cost, apart from $C^{O(h)}$, is at most
\begin{equation}\label{learn:root-saving}
 q^{rv+(h-r)(v-2)}
 =q^{hv-2h+2r}\le q^{-.8h+O(h/\sigma)}.
\end{equation}
This is favorable, and absorbs no attenuation assigned earlier.

\paragraph{Justification of the enumeration.}
The operations above are a way to choose a certificate from a realization,
not an assertion of independence under adaptive choices. Take the union
bound over all resulting static descriptions. More formally, for a fixed
valid tuple, let $I$ be the set of its remaining event requirements and
let $\mathfrak C$ be the family of certificates just described. Write
$\mathcal L(c)$ for the charged direction-batch ledger of $c$, and
$E_{\ell,r}=\{N_r(\ell\cap X_x)>0\}$. The construction proves the
pointwise inequality
\[
 1_{\bigcap_{H\in I}\mathcal B_H}
 \le \sum_{c\in\mathfrak C}
       \prod_{(\ell,r)\in\mathcal L(c)}1_{E_{\ell,r}}.
\]
Every ledger consists of distinct external events, so integration gives
\[
 \Prob\left(\bigcap_{H\in I}\mathcal B_H\,\middle|\,\mathcal O_x\right)
 \le\sum_{c\in\mathfrak C}
       \prod_{(\ell,r)\in\mathcal L(c)}La_\ell.
\]
Different certificates may overlap, which is harmless in this union
bound. Exchange the nonnegative sums over tuples and certificates, and
sum their numerical weights by the encoding order.
Equations~\eqref{learn:overlap-input} and \eqref{learn:line-charge}
give uniform bounds after previously encoded geometry is fixed.
Every encoded hyperplane remains in the fixed typical pencil; in
particular $\lambda'_H<2$ is retained wherever
\eqref{learn:line-charge} is used. The local overlap constraints used
in degree bounds are retained as well. For a residual nonroot, the two
directions remain distinct until its unique plane has been counted;
only their positive weighted sum is then bounded by the product of the
two new-or-old choice bounds $2P+Cp$. Old dictionary references add
multiplicity without adding a hit probability factor.
Discarding other compatibility restrictions can only increase
these weighted sums. This occurs after the component attenuations have been justified
for valid tuples, and never treats an adaptive hit as a fresh event.

There are $O(p)$ dictionary directions and charged events: a pair uses
the fixed number 200, an anchor uses one, and a residual nonroot uses
at most two. Names, traversal orders, component partitions, root choices,
batches, dyads, dictionary references and all polynomial factors in
$p,\sigma,P,R$ contribute at most
$\exp(O(p\log(\sigma+2)))=\exp(o(pP))$.

Finally, a nonsimple component pays for its first two positions by
$e^{-.7P+2\chi}=e^{-.68P}$, and for each additional position by at
most $e^{-P}$. A simple position retains $e^{-.68P}$; its normalized
weighted choice costs at most $e^{2\chi}$ per position apart from the
overhead, or less in \eqref{learn:root-saving}. The deviation factor
$e^\xi$ is also at most $e^{2\chi}$. The weakest per-position saving
comes from sharing $e^{-.68P}$ between the first two positions of a
nonsimple component. It is $e^{-.34P}$ before negligible overhead. Summing
the expansion proves the asserted $B^p e^{-.1pP}$ bound.
\end{proof}

\subsection{Capture, ambient size, and a public-table message}

The two moment bounds now have their promised separate applications.
The second moment controls support points without removing $D_1$;
the high moment is used only for the ambient size.

\begin{lemma}[A successful true-law row]\label{learn:row-success}
For the Poisson batches sampled from $S'$, with probability at least
$c e^{-8P\tau}$ their total size is at most $2qP$ and the test
\eqref{learn:score}, including all sampled points, captures at least
$c n$ points of $S$ and contains at most $n e^{CP}$ ambient points.
The integer $z$ and the score are those determined by these same batches.
\end{lemma}

\begin{proof}
Use the deterministic threshold $z_{\min}/(10q)$. By
Lemma~\ref{learn:high-moment}, the probability that
$|T_x^*|$ exceeds it, conditional on $\mathcal O_x$, is at most
\[
 \exp\{-.1pP+p(b+8P\tau+O(1))\}
 \le e^{-.09pP}\le q^{-50}
\]
outside $D_0\cup D_1\cup D_2$. By
Lemma~\ref{learn:second-moment}, the corresponding bound outside
$D_0\cup D_2$ is $e^{-cP}$. Multiply each conditional estimate by
the probability that $x$ is unsampled and average its own schedules.
These bounds therefore hold for the event that $x$ is unsampled and
its typical contribution exceeds the threshold.

For training points, the expected number of these failures is at most
$n'e^{-cP}$. With failure probability $O(e^{-cP})=o(e^{-8P\tau})$,
there are fewer than $.01n'$ failures outside $D_0\cup D_2$.
For ambient points, the expected number is at most $Q_jq^{-50}$,
so Markov gives at most $n e^{CP}$ failures with a further
$o(e^{-8P\tau})$ exceptional probability. Here $j\le3$ and
$P\tau=o(\sigma)$. The static ambient exceptions $D_1$ are included
by increasing $C$; $D_0\cup D_2$ has size $o(n)$.

Intersect these events with \eqref{learn:empty-good-event}.
Their failure probabilities are subtracted from that event's probability,
not conditioned on its occurrence. The total sample count is Poisson
with mean $qP$, and its exceeding $2qP$ costs $e^{-cqP}$ by
Lemma~\ref{base:tails}, again negligible. The remaining probability
is at least $c e^{-8P\tau}$.

On this joint event, at least $.99n'$ training points satisfy
$z_x\le .1z/q$. Apart from $o(n)$ static exceptions and $.01n'$
unsampled failures, their exceptional and typical score contributions
give $A_x<z/(2q)$, by Lemma~\ref{learn:empty-exceptional} and
$z\ge z_{\min}$. Sampled points are included automatically. This
retains a fixed fraction of $S'$ and hence of $S$.

Conversely, an unsampled point with $z_x\ge .8z/q$, a regular pencil,
and typical contribution of magnitude at most $z_{\min}/(10q)$
has $A_x>z/(2q)$ for sufficiently large $q$. The further exception
depending on $Z$ has size
\[
 O(q^{j+1}/z)\le Cn'e^{b+8P\tau}\le n e^{CP}.
\]
Together with the static exceptions, ambient failures, and at most
$2qP$ sampled points, this proves the size bound after enlarging $C$.
\end{proof}

\begin{lemma}[Implementation by proposals]\label{learn:implementation}
The true-law success in Lemma~\ref{learn:row-success} gives the
description scheme of Proposition~\ref{learn:short-row}.
\end{lemma}

\begin{proof}
Transmit the prepared plane list and cell sizes, the integers $n,n'$,
and the choice of the present mode. These cost $O(q+\sigma)$ bits.
The decoder can now generate a public-table proposal consisting of
$R$ independent Poisson batch sizes of mean $qL$, with entries uniform
in its known domain $U_S$. The table supplies independent repetitions
of this entire proposal. It is fresh relative to the data used to
prepare $S'$.

The encoder accepts a proposal precisely if every entry belongs to
$S'$ and the conclusion of Lemma~\ref{learn:row-success} holds for
that row. The encoder knows the hidden data, so this is a well-defined
finite test. For a specified ordered row with total count $m$ and all
entries in $S'$, the proposal-to-true-law probability ratio is
\[
 \left(\frac{n'}{|U_S|}\right)^m.
\]
The Poisson probabilities of the batch sizes cancel exactly. On
successful true-law rows $m\le2qP$ and $n'\ge n/2$, so summing these
ratios over the successful rows proves proposal success probability
at least
\[
 c e^{-8P\tau}e^{-2qP(d_S+\log2)}
 \ge e^{-C_1qP(d_S+1)}.
\]
Choose a fixed $C_2>C_1$ and search at most
\[
 M=\left\lceil e^{C_2qP(d_S+1)}\right\rceil
\]
independent proposals, taking the first success and declaring failure
if none exists. The failure probability is at most
\[
 \exp\{-M e^{-C_1qP(d_S+1)}\}
 \le \exp\{-e^{(C_2-C_1)qP(d_S+1)}\}
 \le C e^{-cq}.
\]
The successful index has a self-delimiting description of
$O(qP(d_S+1))$ bits. The encoder also sends $z\le Q_j$, costing
$O(\sigma)$ bits. The decoder reconstructs that proposal, computes
the score from its samples and the transmitted list data, includes
the sampled points, and intersects the resulting set with $U_S$.
This intersection preserves all captured points of $S'$.

The decoder neither repeats the acceptance test nor needs the hidden
supports to evaluate the resulting score. The guarantee concerns the
first accepted row within a finite cutoff; it does not assert that
this accepted row has an unbiased true-law distribution. Any later
validation procedure uses its own fresh table. Finally, conditional
on arbitrary earlier readiness data, the fixed-input calculation above
is unchanged. This proves the uniformity, length, and failure claims.
\end{proof}

\begin{proof}[Completion of Proposition~\ref{learn:short-row}]
Lemmas~\ref{learn:row-success} and~\ref{learn:implementation} give
the stated set, message bound, and failure probability.
\end{proof}

\begin{proof}[Completion of the description lemma]
\label{learn:completion}
The only unresolved branch in the proof of
Lemma~\ref{geom:description} was the prepared small-cell case in
projective dimension two or three. Its inputs satisfy exactly the
hypotheses of Proposition~\ref{learn:short-row}, which supplies that
branch at a cost bounded by the lemma's stated budget. Dimension
two has no plane-list reduction and no external radial overlaps.
The dimension-three large-cell reduction of
Proposition~\ref{geom:preparation} invokes the now established
dimension-two result. The bounded higher-dimensional projection
induction of Proposition~\ref{geom:projection} then applies successively.
Its validation calls use fresh
tables, as stipulated in the preceding section. Hence all branches
of Lemma~\ref{geom:description} are proved, with the specified
one-sided capture, message length, and exponentially small failure.
\end{proof}
